\documentclass[11pt]{article}

\usepackage[T1]{fontenc}
\usepackage{lmodern}
\usepackage{amsmath,amssymb,amsthm}
\usepackage[margin=1in]{geometry}
\usepackage{enumitem}
\usepackage{xcolor}
\usepackage{microtype}
\usepackage[hidelinks]{hyperref}
\hypersetup{
  pdftitle={Paper 1, Theorems 1.2--1.3: Statement Amendment},
  pdfauthor={ShenWork formalization project}
}

\newcommand{\R}{\mathbb{R}}
\newcommand{\abs}[1]{\left\lvert #1\right\rvert}
\setlength{\emergencystretch}{2em}

\title{Paper 1, Theorems 1.2--1.3: Statement Amendment}
\author{Internal research note for the \texttt{ShenWork} formalization}
\date{14 July 2026}

\begin{document}
\maketitle

\begin{abstract}
The stability proof in Section~5 of Shen's traveling-wave paper has seven
independent statement-level defects.  The perturbed quadratic in (5.31) is
positive at the unperturbed tail exponent, so the root comparison in (5.35)
has the wrong direction; the weighted norm in (1.21) is written in laboratory
rather than moving coordinates; and the exponential factor following (5.35)
has the wrong sign.  In addition, the chemotactic flux difference in (5.18)
has two reversed signs.  Finally, (5.27) loses one factor of $b_1$ in
Young's inequality and (5.29)--(5.30) drop the
$M^{2(\gamma-1)}$ factor from Lemma~5.3.  Finally, Lemma~5.2 is
one-sided in $U_x/U$, whereas (5.23) uses it as an absolute bound.  This note gives the direct amendment supported by
the corrected calculation, explains the corresponding change to the
uniqueness theorem, and records the exact Lean certificates and remaining
analytic boundary.
\end{abstract}

\section{Source checked}

The source examined is the original submission, not a repository paraphrase:
\begin{quote}
Wenxian Shen, \emph{Existence, uniqueness, stability, and monotonicity of
traveling waves for repulsion/attraction chemotaxis models with logistic type
source}, arXiv:2605.04401v1, 6 May 2026.
\end{quote}
The archived local copy is \path{paper1.pdf}.  Theorem~1.2 claims weighted
stability for every
\[
  \kappa<\eta<\frac{1}{1+\abs{\chi}^{1/6}},
\]
where $\kappa=(c-\sqrt{c^{2}-4})/2$, and writes the weighted conclusion in
laboratory coordinates.

\section{The root comparison in (5.35)}

The energy coefficient obtained in (5.31) is
\begin{equation}
 q(\eta)=\eta^{2}-(c-A)\eta+(1+B),
 \label{eq:q}
\end{equation}
where, in the notation of (5.32)--(5.33),
\[
 A=\abs{\chi}^{1-3\sigma}D',
 \qquad
 B=\abs{\chi}^{1-3\sigma}D''.
\]
The displayed formulas give $A\geq0$ and, when $\chi\neq0$, $B>0$.
The unperturbed tail exponent satisfies
\[
  \kappa^{2}-c\kappa+1=0.
\]
Substitution into \eqref{eq:q} therefore gives the exact identity
\begin{equation}
  q(\kappa)=A\kappa+B>0.
  \label{eq:q-kappa}
\end{equation}
Consequently $\kappa$ cannot lie between the two roots of $q$.  Under the
paper's speed inequality it lies strictly below the perturbed lower root:
\[
  \kappa<\kappa^{-}_{c,\chi,m,\alpha,\gamma}.
\]
By continuity, \eqref{eq:q-kappa} also produces weights strictly above
$\kappa$ for which $q(\eta)>0$.  Thus excluding only the endpoint
$\eta=\kappa$ does not repair (5.35).  The global energy estimate proves
decay only above the perturbed lower root.

\section{The coordinate mismatch in (1.21)}

Section~5 changes to the moving coordinate $z=x-ct$ and estimates
\begin{equation}
 E_{\rm move}(t)=\int_{\R}e^{2\eta z}
   \abs{u(t,z+ct)-U(z)}^{2}\,dz.
 \label{eq:move}
\end{equation}
The printed formula (1.21) instead uses
\begin{equation}
 E_{\rm lab}(t)=\int_{\R}e^{2\eta x}
   \abs{u(t,x)-U(x-ct)}^{2}\,dx.
 \label{eq:lab}
\end{equation}
The substitution $x=z+ct$ gives
\begin{equation}
 E_{\rm lab}(t)=e^{2\eta ct}E_{\rm move}(t).
 \label{eq:change}
\end{equation}
Hence decay of \eqref{eq:move} does not imply decay of \eqref{eq:lab} without
an additional rate stronger than $2\eta c$, which the proof does not supply.
The natural correction is to use \eqref{eq:move}, equivalently the weight
$e^{2\eta(x-ct)}$ in laboratory coordinates.

\section{The sign after (5.35)}

The paper defines the value $\lambda=q(\eta)$ to be negative and then writes
an upper bound with $e^{-\lambda t}$, claiming that it tends to zero.  For
$\lambda<0$, however, $-\lambda>0$ and this factor grows.  With the displayed
definition the decaying factor is $e^{\lambda t}$.  Equivalently, define the
positive rate $\rho=-\lambda$ and write $e^{-\rho t}$.

\section{The two chemotactic sign errors in (5.18)}

The elliptic equation in (1.1) is
\[
 v_{xx}-v+u^\gamma=0,
 \qquad\text{hence}\qquad v_{xx}=v-u^\gamma.
\]
Put $w=u-U$ and $z=v-V$.  The zero-order part of the difference of the
two chemotactic flux derivatives is
\begin{equation}
 \bigl(v a_m-a_{m+\gamma}\bigr)w+U^m z.
 \label{eq:correct-flux-zero}
\end{equation}
After multiplication by $-\chi$, the perturbation equation must therefore
contain
\begin{equation}
 -\chi\bigl(v a_m-a_{m+\gamma}+b_2\bigr)w-\chi b_4 z.
 \label{eq:correct-perturbation-zero}
\end{equation}
Equation (5.18) instead prints $+a_{m+\gamma}$ inside the first parentheses
and $+\chi b_4z$; the same two reversals propagate into (5.19).

The conservative energy budget can nevertheless be repaired without a new
asymptotic constant.  For either sign of $\chi$,
\[
 \abs{v a_m-a_{m+\gamma}}
 \leq (2m+\gamma)M^{m+\gamma-1},
\]
and the $b_4z$ term is estimated by absolute value.  The
$(2m+\gamma)M^{m+\gamma-1}$ contribution is already present in (5.33).

\section{The lost square in (5.27)}

The $J_1$ cross term obeys the direct Young estimate
\[
 \abs{\chi b_1 W_xW}
 \leq \frac12 W_x^2+\frac12\abs{\chi}^2 b_1^2W^2.
\]
Formula (5.27), propagated into (5.33), uses $b_1$ rather than $b_1^2$.
There is no hypothesis making this replacement an upper bound; in
particular, the negative-$\chi$ branch allows arbitrary $\abs{\chi}$.
If $B_1$ bounds $\abs{b_1}$, the corrected constant contribution is
$\abs{\chi}^2B_1^2/2$.

\section{The missing resolver power in (5.29)--(5.30)}

Lemma~5.3 gives
\[
 \int_{\R}V^2\leq
 \frac{\gamma^2M^{2(\gamma-1)}}{(1-\eta)^2}\int_{\R}W^2,
 \qquad
 \int_{\R}V_x^2\leq
 \frac{\gamma^2M^{2(\gamma-1)}}{1-\eta^2}\int_{\R}W^2.
\]
The subsequent estimates retain $\gamma^2$ but omit
$M^{2(\gamma-1)}$.  This power is not identically one: the positive-$\chi$
asymptotic cap $M_\chi$ can exceed one.  For $\chi\neq0$, put
\[
 x=\abs{\chi}^{1/6},\qquad
 K=1+\gamma^2M^{2(\gamma-1)}\frac{(1+x)^2}{x^2}.
\]
Whenever $0\leq\eta<1/(1+x)$, this $K$ bounds both $1+R_V(\eta)$
and $1+R_{V_x}(\eta)$.  The case $\chi=0$ is separate and all
chemotactic corrections vanish.

\section{The one-sided logarithmic-derivative bound used in (5.23)}

Lemma~5.2 proves only
\[
  \frac{U_x(x)}{U(x)}\leq C.
\]
Case ii.3 then estimates $\abs{b_2}$ by using the stronger statement
$\abs{U_x/U}\leq C$.  This implication is false.  For a decreasing wave,
$U_x/U\leq0$ makes the conclusion of Lemma~5.2 automatic, while its negative
part is precisely what the absolute value retains.

There is a direct repair.  Strengthen the speed threshold so that the
logarithmic Riccati vector field points inward at both $-1$ and $1$.
Two applications of the right-derivative fencing theorem, starting from
$U_x/U\to0$ at the left end, then give
\[
  -1\leq \frac{U_x}{U}\leq1,
  \qquad \text{and hence}\qquad
  \abs{U_x/U}\leq1.
\]
This bound is independent of the wave speed and requires no monotonicity, so
it applies to both sensitivity signs.  Without the stronger barrier-speed
condition, an absolute logarithmic-derivative bound must instead remain an
explicit input to the $B_2$ budget; the one-sided Lemma~5.2 does not provide
it.

\section{Recommended amended statement}

Let $B_i$ bound $\abs{b_i}$ and use the corrected common resolver factor
$K$ above.  The valid replacements for (5.32)--(5.33) are
\begin{align*}
 A_\chi&=\abs{\chi}B_1+\frac{\abs{\chi}B_3}{2}K,\\
 B_\chi&=\abs{\chi}\bigl((2m+\gamma)M^{m+\gamma-1}+B_2\bigr)
   +\frac{\abs{\chi}^2B_1^2}{2}
   +\frac{\abs{\chi}(B_3+B_4)}{2}K.
\end{align*}
Here $B_2$ must be supplied by an absolute logarithmic-derivative estimate.
At the stronger two-sided Riccati barrier speed, any wave may use the fixed
bound above.  Without that condition, the genuine speed dependence of an
absolute bound must be retained.
Set
\begin{align*}
 \Delta&=(c-A_{\chi})^{2}-4(1+B_{\chi}),\\
 \kappa^-&=\frac{(c-A_{\chi})-\sqrt{\Delta}}{2},\\
 \kappa^+&=\frac{(c-A_{\chi})+\sqrt{\Delta}}{2}.
\end{align*}
Choose the speed threshold so that $\Delta>0$ and
\[
 \kappa^-<\frac{1}{1+\abs{\chi}^{1/6}}<\kappa^+.
\]

\medskip
\noindent
\fbox{%
\begin{minipage}{0.92\textwidth}
\textbf{Recommended amendment to Theorem 1.2.}
Replace the condition $\kappa<\eta$ by
\[
 \kappa^-_{c,\chi,m,\alpha,\gamma}<\eta
   <\frac{1}{1+\abs{\chi}^{1/6}},
\]
and replace (1.21) by
\[
 \lim_{t\to\infty}\int_{\R}e^{2\eta z}
   \abs{u(t,z+ct;u_{0})-U^{*}(z)}^{2}\,dz=0.
\]
Retain the uniform conclusion (1.22), subject to the compactness and
far-left rigidity argument in Step~4.  Write the energy decay as
$e^{\lambda t}$ for $\lambda<0$, or as $e^{-\rho t}$ for
$\rho=-\lambda>0$.
\end{minipage}}

\medskip
Theorem~1.3 uses Theorem~1.2 with the second wave as initial data.  Its common
tail exponent must therefore leave room for the corrected stability weight:
the hypothesis should require $\kappa_{1}>\kappa^-$, not merely
$\kappa_{1}>\kappa$.  The waves constructed in Theorem~1.1 satisfy a family
of tail estimates, so such a rate can be selected once $\kappa^-$ lies below
the construction's tail cap.

\section{What would be needed to retain the original interval}

Equation \eqref{eq:q-kappa} obstructs only the \emph{global coefficient}
estimate used in (5.31); it is not a counterexample to stability itself.
The perturbation coefficients may decay in the right tail.  Retaining every
$\eta>\kappa$ would require a new spatially localized coercivity or weighted
spectral argument.  No such argument appears in Section~5.  It would be a new
theorem, not a formalization of the written proof.

\section{Lean certificates and current boundary}

The repository contains zero-\texttt{sorry}, standard-axiom certificates.
The algebraic root certificates in
\path{ShenWork/Paper1/Theorem12RootObstruction.lean} are
\begin{itemize}[leftmargin=2em]
\item \path{paper531_kappa_not_between_perturbed_roots};
\item \path{paper531_kappa_lt_rootMinus};
\item \path{paper531_positive_inside_stated_weight_window}.
\end{itemize}
The coordinate certificate
\path{laboratoryWeightedL2Energy_eq_exp_mul_coMoving}, together with a
non-vacuous coordinate witness, is in
\path{ShenWork/Paper1/Theorem12CoordinateAudit.lean}.  The two sign
certificates
\path{paper531_printed_decay_factor_tendsto_atTop} and
\path{paper531_corrected_decay_factor_tendsto_zero} are in
\path{ShenWork/Paper1/Theorem12RootObstruction.lean}.
The corrected flux identity
\path{paper5ChemFluxDifference_expansion_corrected} and its uniform
zero-order estimate
\path{paper5CorrectedChemZeroCoefficient_abs_le} are in
\path{ShenWork/Paper1/Theorem12MeanCoefficients.lean}.
The corrected Young, resolver-cap, and scalar-budget certificates
\path{paper5J1VariableDensity_le},
\path{paper5WeightedResolverFactors_le_cap}, and
\path{paper5ResolvedW2Coefficient_le_corrected531} are in
\path{ShenWork/Paper1/Theorem12WeightedEnergy.lean}.
The fencing and absolute logarithmic-derivative certificates
\path{lower_barrier_of_tendsto_atBot} and
\path{abs_waveLogDerivative_le_one_of_barrier_speed} are in
\path{ShenWork/Paper1/Theorem12LogDerivative.lean}.  Their fixed coefficient
bundle \path{paper5CoefficientBounds_of_barrier_speed_corrected_wave} is in
\path{ShenWork/Paper1/Theorem12MeanCoefficients.lean}.

\path{ShenWork/Paper1/Theorem12Corrected.lean} proves the scalar Gronwall and
moving-frame localization wiring, but keeps the general-data whole-line PDE
energy and compactness input explicit.  The full amended stability theorem is
therefore not yet claimed as proved.  The corrected structural reduction for
Theorem~1.3 is in
\path{ShenWork/Paper1/Theorem13Corrected.lean}; its remaining analytic input
is the same bounded-solution stability theorem.

\end{document}
